\clearpage\begingroup
\SourceLeaf{176}\NoSourceHead
\begin{center}
\textbf{\Large VERMISCHTES ZUR PHYSIK.}\par\bigskip
\textbf{\large [ZUR AKUSTIK.]}\par\bigskip
\textbf{\large\textls[120]{BRIEFWECHSEL.}}\par\bigskip
\textbf{\textsc{Gauss} an \textsc{W. Weber}.}
\end{center}
\noindent\hspace*{2\parindent}Wohlgeborner Herr\par
\noindent\hspace*{2\parindent}Hochzuverehrender Herr Professor

Ihr gütiges Schreiben nebst Ihren schätzbaren akustischen Abhandlungen erhielt ich zu einer Zeit, wo die Überhäufung mit Geschäften mich nöthigte, letztere einstweilen bei Seite zu legen, und mir ihre Lectüre auf eine etwas freiere Zeit vorzubehalten. Gegenwärtig sind wenigstens die Ferien — oder in Beziehung auf mich richtiger, ist die Zeit der Ferien — eingetreten. Ich habe Ihre Aufsätze zur Hand genommen; allein wie lebhaft mich auch die Originalität Ihrer Arbeiten interessirt hat, bin ich doch bald gewahr geworden, dass sie zu einer gründlichen Beurtheilung des reichhaltigen Inhalts durchaus eine nachhaltigere Beschäftigung erfordern und verdienen, als mir gegenwärtig vergönnt ist, ihnen widmen zu können. Ich darf es aber nicht länger verschieben, Ihnen für die gewogentliche Mittheilung meinen ergebensten Dank und zugleich meine Freude zu bezeugen, dass Sie sich diesen interessanten Untersuchungen mit so viel Eifer als Erfolg widmen.

Ich bin immer der Meinung gewesen, dass die Akustik zu denjenigen Theilen der mathematischen Physik gehört, wo noch die allerglänzendsten
\par\HistoricalFoot{}{22\textsuperscript{*}}
\clearpage\SourceLeaf{177}\HistoricalHead{172}{L}{PHYSIK.\quad BRIEFWECHSEL.}
\noindent Fortschritte zu machen sind. In der That handelt es sich dabei bloss um räumliche und Zeit-Relationen und der Gegenstand sollte also ganz der Mathematik unterwürfig gemacht werden können; und doch wie wenig, wie äusserst wenig wissen wir noch! Bei den dem Anscheine nach gemeinsten Dingen wissen wir noch gar nicht einmal, wie wir sie angreifen sollen. Alle unsere bisherigen Untersuchungen beschränken sich 1) auf die Geschwindigkeit der Fortpflanzung und 2) auf die Dicke der Schallwellen; von dem Verhältniss dieser beiden Grössen hängt die Höhe der Töne ab. Aber das Unterscheidende, Eigenthümliche der Töne, was in einigen Fällen durch das Wort timbre bezeichnet wird, aber am wunderbarsten in den sogenannten articulirten Tönen sich äussert, dies ist, insofern mathematisch klare Einsicht gefordert wird, bis jetzt eine völlige terra incognita. Es scheint, dass dies nur von der Gestalt der Schallwellen abhängen kann. Welch ein unermessliches Feld liegt hier noch vor uns. Ich bin überzeugt, dass der menschliche Geist es einst öffnen, und hier dieselbe Klarheit schaffen wird, die den optischen Wissenschaften gegeben ist. Man sollte selbst glauben, dass für den, welcher lebendigen Eifer, mathematische Kraft, Anstelligkeit zu Versuchen und hinreichende Musse vereinigte, es so gar schwer nicht werden könnte, in diese Dunkelheit das erste Licht zu bringen. Es handelt sich nur darum, zu lehren, welches die specifischen Unterschiede für verschiedene Töne sind, nicht im schallgebenden Körper und nicht im Ohr, sondern in dem elastischen Medium zwischen beiden. In Beziehung auf den schallgebenden Körper ist für diesen Gegenstand durch \textsc{von Kempelens} genialisches Werk[*)]\begingroup\renewcommand{\thefootnote}{}\footnotetext{[*) \textsc{Wolfgang von Kempelen}, \textit{Mechanismus der menschlichen Sprache}, Wien 1791.]}\endgroup\ schon so sehr viel vorgearbeitet, dass man sich wundern muss, dass wir seitdem (seit 40 Jahren) noch nichts weiter gekommen sind. Herzlich wünsche ich, dass es Ihnen vorbehalten sein möge, hierin eine neue Bahn zu brechen.

Unter Versicherung meiner aufrichtigsten Hochachtung empfehle ich mich Ihrem freundschaftlichen Andenken
\par\medskip
\noindent\hspace*{\parindent}Göttingen den 2. April 1830.\hfill\begin{tabular}[t]{@{}c@{}}ganz ergebenst\\C. F. \textsc{Gauss}.\end{tabular}\hspace*{\parindent}
\par\vspace{2\baselineskip}\begin{center}\rule{.25\linewidth}{.4pt}\end{center}
\clearpage\endgroup
